\documentclass{article}
\usepackage{amsmath, amssymb}
\usepackage{ctex}
\usepackage{geometry}
\geometry{a4paper, margin=1in}

\begin{document}
\title{第七章 机械振动 例题答案}
\maketitle

\begin{enumerate}
\item 设物体的运动学方程为 $x=A\sin(\omega t+\phi)$。

由 $A=0.08\,\mathrm{m}$，$T=4\,\mathrm{s}$，得 $\omega=\dfrac{2\pi}{T}=\dfrac{\pi}{2}\,\mathrm{rad/s}$。

代入 $t=0$ 时 $x=0.04\,\mathrm{m}$，得 $0.04=0.08\sin\phi$，即 $\sin\phi=\dfrac{1}{2}$，故 $\phi=\dfrac{\pi}{6}$ 或 $\dfrac{5\pi}{6}$。由 $v_0>0$ 得 $\phi=\dfrac{\pi}{6}$。

故运动学方程为 $x=0.08\sin\left(\dfrac{\pi}{2}t+\dfrac{\pi}{6}\right)\,\mathrm{m}$，速度为 $v=0.08\times\dfrac{\pi}{2}\cos\left(\dfrac{\pi}{2}t+\dfrac{\pi}{6}\right)\,\mathrm{m/s}$。

当 $t=1.0\,\mathrm{s}$ 时：
\begin{align*}
x &= 0.08\sin\left(\dfrac{\pi}{2}+\dfrac{\pi}{6}\right) = 0.08\sin\dfrac{2\pi}{3} = 0.08\times\dfrac{\sqrt{3}}{2} \approx 0.07\,\mathrm{m},\\
v &= 0.08\times\dfrac{\pi}{2}\cos\left(\dfrac{\pi}{2}+\dfrac{\pi}{6}\right) = 0.04\pi\times\left(-\dfrac{1}{2}\right) \approx -0.06\,\mathrm{m/s},\\
f &= ma = -0.01\times 0.08\times\left(\dfrac{\pi}{2}\right)^2\sin\left(\dfrac{\pi}{2}+\dfrac{\pi}{6}\right) \approx -1.71\times 10^{-3}\,\mathrm{N}.
\end{align*}

\item 设运动学方程为 $x=A\sin(\omega t+\phi)$，其中 $A=0.08\,\mathrm{m}$，$\omega=\dfrac{\pi}{2}\,\mathrm{rad/s}$。

当 $t=0$ 时，$x=0$，代入得 $0=0.08\sin\phi$，故 $\phi=0$ 或 $\pi$。由 $t=0$ 时 $v_0<0$，只能取 $\phi=\pi$。

故运动学方程为 $x=0.08\sin\left(\dfrac{\pi}{2}t+\pi\right)\,\mathrm{m}$。

设由起始 $x=0$ 到 $x=-8\,\mathrm{cm}$ 所需的最短时间为 $t'$，则
\[
-0.08 = 0.08\sin\left(\dfrac{\pi}{2}t'+\pi\right),\quad \sin\left(\dfrac{\pi}{2}t'+\pi\right)=-1,
\]
即 $\dfrac{\pi}{2}t'+\pi=\dfrac{3\pi}{2}$，解得 $t'=1\,\mathrm{s}$。

\item 设运动学方程为 $x=A\cos(\omega t+\phi)$，$A=6\,\mathrm{cm}$。

由 $t=0$ 时 $x_0=-3\,\mathrm{cm}$，得 $-3=6\cos\phi$，即 $\cos\phi=-\dfrac{1}{2}$，故 $\phi=\dfrac{2\pi}{3}$ 或 $\dfrac{4\pi}{3}$。由 $v_0>0$ 取 $\phi=\dfrac{4\pi}{3}$。

再由 $t=1.5\,\mathrm{s}$ 时 $x=3\,\mathrm{cm}$，得 $3=6\cos\left(\omega\times 1.5+\dfrac{4\pi}{3}\right)$，即 $\omega\times 1.5+\dfrac{4\pi}{3}=\dfrac{5\pi}{3}$（由 $v>0$ 取此支）。

解得 $\omega=\dfrac{2\pi}{9}\,\mathrm{rad/s}$。

故振动方程为 $x=6\cos\left(\dfrac{2\pi}{9}t-\dfrac{4\pi}{3}\right)\,\mathrm{cm}$。

\item \textbf{(1)} 选平衡位置为坐标原点，向下为 $x$ 轴正方向。

平衡时 $mg=kl$，故 $k=\dfrac{mg}{l}$。

当位移为 $x$ 时，由牛顿第二定律 $m\dfrac{d^2x}{dt^2}=mg-k(l+x)=-kx$，
即 $\dfrac{d^2x}{dt^2}+\dfrac{k}{m}x=0$。

故物体作简谐振动，角频率 $\omega=\sqrt{\dfrac{k}{m}}=\sqrt{\dfrac{g}{l}}=\sqrt{\dfrac{9.8}{0.098}}=10\,\mathrm{rad/s}$。

由初始条件 $x_0=0$，$v_0=1\,\mathrm{m/s}$，得
\[
A=\sqrt{x_0^2+\dfrac{v_0^2}{\omega^2}}=\sqrt{0+\dfrac{1}{100}}=0.1\,\mathrm{m},
\]
$\tan\phi=-\dfrac{v_0}{\omega x_0}\to-\infty$，即 $\phi=-\dfrac{\pi}{2}$。

故振动方程为 $x=0.1\cos\left(10t-\dfrac{\pi}{2}\right)\,\mathrm{m}$（SI）。

\textbf{(2)} 最大回复力 $F_{\max}=kA=\dfrac{mg}{l}\cdot A=\dfrac{0.1\times 9.8}{0.098}\times 0.1=1\,\mathrm{N}$。

\item 系统受力矩 $M=-mgl\sin\theta$。当 $\theta$ 很小时 $\sin\theta\approx\theta$，故 $M\approx -mgl\theta$。

由转动定律 $M=J\dfrac{d^2\theta}{dt^2}$，其中单摆 $J=ml^2$，得
\[
ml^2\dfrac{d^2\theta}{dt^2}=-mgl\theta,\quad \text{即}\quad \dfrac{d^2\theta}{dt^2}+\dfrac{g}{l}\theta=0.
\]

令 $\omega=\sqrt{\dfrac{g}{l}}$，则 $\dfrac{d^2\theta}{dt^2}+\omega^2\theta=0$，故单摆作简谐振动。

周期 $T=\dfrac{2\pi}{\omega}=2\pi\sqrt{\dfrac{l}{g}}$。

\item 在平衡位置处 $mg=kl$，故 $k=\dfrac{mg}{l}$。

以平衡位置为坐标原点，向下为 $x$ 轴正方向。当物体运动到 $x$ 处时，对物体应用牛顿第二定律，对滑轮应用转动定律：
\begin{align*}
mg - T_1 &= m\dfrac{d^2x}{dt^2},\\
(T_1 - T_2)R &= \dfrac{1}{2}MR^2\alpha,\\
T_2 &= k(x+l),\\
\alpha &= \dfrac{1}{R}\dfrac{d^2x}{dt^2}.
\end{align*}

联立以上方程，消去 $T_1$、$T_2$ 和 $\alpha$，利用平衡条件 $mg=kl$ 化简得
\[
\dfrac{d^2x}{dt^2}+\dfrac{k}{m+M/2}x=0.
\]

故系统作简谐振动，振动圆频率 $\omega=\sqrt{\dfrac{k}{m+M/2}}$。

\item 平衡时弹簧伸长 $x_0$，由合外力矩为零得 $kx_0 l=mg\dfrac{l}{2}$，即 $x_0=\dfrac{mg}{2k}$。

当杆摆到角度 $\theta$ 时，弹簧伸长量为 $x_0+x$（其中 $x\approx l\theta$）。合外力矩为
\[
M = mg\dfrac{l}{2}\cos\theta - 2mgl\sin\theta - k(x_0+x)l\cos\theta.
\]

小角度下 $x\approx l\theta$，$\cos\theta\approx 1$，$\sin\theta\approx\theta$，代入并利用平衡条件化简得
\[
M=-(2mgl+kl^2)\theta.
\]

由刚体定轴转动定律 $J\dfrac{d^2\theta}{dt^2}=M$，且 $J=\dfrac{1}{3}ml^2+\dfrac{1}{3}(2m)(2l)^2=3ml^2$，得
\[
\dfrac{d^2\theta}{dt^2}+\dfrac{2mgl+kl^2}{3ml^2}\theta=0.
\]

故杆作简谐振动，角频率 $\omega=\sqrt{\dfrac{2mg+kl}{3ml}}$，周期 $T=2\pi\sqrt{\dfrac{3ml}{2mg+kl}}$。

\item \textbf{(1)} $t=0$ 时音叉尖端在平衡位置 $x=0$ 且 $v<0$，由旋转矢量图得初相 $\varphi=\dfrac{\pi}{2}$。

故运动学方程为 $x=0.1\cos\left(6.28\times 10^2 t+\dfrac{\pi}{2}\right)\,\mathrm{cm}$。

\textbf{(2)} $t=0$ 时音叉尖端 $x=-A/2$ 且 $v<0$，由旋转矢量图得初相 $\varphi=\dfrac{2\pi}{3}$。

故运动学方程为 $x=0.1\cos\left(6.28\times 10^2 t+\dfrac{2\pi}{3}\right)\,\mathrm{cm}$。

\item 在液体平衡时的液面上取一点 $O$ 为 $x$ 轴原点，$x$ 轴向上为正。时刻 $t$ 液体的动能
\[
E_k=\dfrac{1}{2}Sl\rho \dot{x}^2.
\]

只需计算 $x$ 小段液柱从左边移到右边所需做的功即得重力势能变化量：
\[
E_p=\rho S x^2 g.
\]

根据机械能守恒定律：
\[
E=E_k+E_p=\dfrac{1}{2}Sl\rho\dot{x}^2+\rho Sx^2 g=\text{常量}.
\]

对 $t$ 求导并整理得
\[
\ddot{x}+\dfrac{2g}{l}x=0.
\]

故液柱作简谐振动，振动周期 $T=2\pi\sqrt{\dfrac{l}{2g}}$。

\end{enumerate}

\end{document}
